\documentclass[11pt]{article}
\usepackage[margin=1in]{geometry}
\usepackage{amsmath,amssymb,amsthm,mathtools,booktabs}
\usepackage{hyperref}
\hypersetup{colorlinks=true,linkcolor=blue,citecolor=blue,urlcolor=blue,
pdftitle={Full asymptotic expansions for partition sums of acyclic orientations}}
\setlength{\parskip}{.35em}
\setlength{\parindent}{1.1em}
\setlength{\emergencystretch}{2em}
\newtheorem{theorem}{Theorem}[section]
\newtheorem{lemma}[theorem]{Lemma}
\newtheorem{corollary}[theorem]{Corollary}
\newtheorem{proposition}[theorem]{Proposition}
\theoremstyle{remark}\newtheorem{remark}[theorem]{Remark}
\newcommand{\AO}{\operatorname{AO}}
\newcommand{\E}{\mathbb E}
\newcommand{\eps}{\varepsilon}
\newcommand{\ii}{\mathrm i}
\newcommand{\dd}{\,\mathrm d}
\newcommand{\coeff}[2]{[#1^{#2}]}
\title{Full asymptotic expansions for partition sums of acyclic orientations}
\author{}
\date{1 October 2026}
\begin{document}
\maketitle
\begin{abstract}
We obtain the full prefactors and arbitrary finite orders of the asymptotic
expansions for OEIS A372395 and A370613, which sum acyclic orientations of
complete multipartite graphs over unrestricted and distinct-part integer
partitions. The normalized counts have the respective forms
$A_+n^{-1}e^{C_+\sqrt n}$ and $A_-n^{-3/4}e^{C_-\sqrt n}$, with corrections
in powers of $n^{-1/2}$. The constants $C_\pm$ were previously identified
by Sun. The amplitudes include a nontrivial contribution from a tilted
Gamma integral. A uniform quadratic majorant supplies relative-precision
tail control despite the signed polynomial integral representation.
We give a finite coefficient algorithm, explicit first and second
corrections, reproducible exact checks through $n=300$, and controlled
Lambert-$W$ inverse expansions. The result is an all-orders Poincar\'e
expansion; it does not assert an exponentially complete transseries.
\end{abstract}

\section{Results and relation to prior work}

For an integer partition $\lambda$ of $n$, let $K_\lambda$ be the complete
multipartite graph with part sizes $\lambda_1,\lambda_2,\ldots$. Define
\[
 B_+(n)=\sum_{\lambda\vdash n}\AO(K_\lambda),\qquad
 B_-(n)=\sum_{\substack{\lambda\vdash n\\\lambda\text{ distinct}}}\AO(K_\lambda).
\]
These are OEIS A372395 and A370613 \cite{oeisb,oeisf}. Each partition
profile occurs once; there is no additional multinomial choice of vertex
sets. We set both empty counts equal to one.

Sun \cite{sun} derives the exact Gamma representation used below and proves
$\log(B_\pm(n)/n!)\sim C_\pm\sqrt n$. His Section 11 explicitly asks for
the prefactors and complete expansions. The new refinement here is at
relative precision: the amplitudes, every algebraic correction, and the
associated inverse. The main additional estimate is a uniform quadratic
majorant which removes the large-part cutoff without losing an
$e^{o(\sqrt n)}$ factor.

Use $\eta=+1$ for unrestricted partitions and $\eta=-1$ for distinct parts.
For real $c$ in its natural domain, put
\begin{align}
 \rho_\eta(x;c)&=\frac1{e^{cx+x^2/2}-\eta},\label{eq:rho}\\
 J_\eta(c)&=-\eta\int_0^\infty
       \log(1-\eta e^{-cx-x^2/2})\dd x,\\
 M_j(c)&=\int_0^\infty x^j\rho_\eta(x;c)\dd x,\qquad j\ge1.
\end{align}
The Bose case requires $c>0$; the Fermi case allows every real $c$.
There is a unique $c=c_\eta$ such that $M_1(c)=1$. Indeed $M_1$ is
continuous and strictly decreasing, with limits infinity and zero at the
two endpoints of its domain. In the Bose case the divergence as
$c\downarrow0$ is logarithmic at the origin; in the Fermi case it follows
from the interval where $cx+x^2/2<0$ as $c\to-\infty$.

All unlabelled $M_j,J$ below are evaluated at this saddle unless their
argument is shown. Set
\begin{gather}
 C=c+J(c),\qquad
 V=J''(c)=\int_0^\infty x^2\rho(1+\eta\rho)\dd x>0,\label{eq:CV}\\
 E_0=\exp\left(\frac12-\frac{M_3}{3}+\frac{M_2^2}{8}\right),\label{eq:E0}\\
 p_+=1,\quad A_+=\frac{\sqrt{c_+}\,E_{0,+}}{2\pi\sqrt{V_+}},
 \qquad
 p_-=\frac34,\quad A_-=\frac{E_{0,-}}{2\sqrt{\pi V_-}}.\label{eq:A}
\end{gather}

\begin{theorem}\label{thm:main}
For each $\eta\in\{+1,-1\}$ there are real constants $a_{\eta,j}$,
with $a_{\eta,0}=1$, such that for every fixed integer $L\ge0$,
\begin{equation}\label{eq:main}
 B_\eta(n)=A_\eta n!n^{-p_\eta}e^{C_\eta\sqrt n}
 \left(\sum_{j=0}^{L}a_{\eta,j}n^{-j/2}
             +O_L(n^{-(L+1)/2})\right).
\end{equation}
Section~\ref{sec:generator} gives a finite algorithm for every coefficient.
\end{theorem}

\begin{center}\small
\begin{tabular}{@{}lrr@{}}
\toprule
 & Unrestricted & Distinct parts\\
\midrule
$c_\eta$ & $0.764996442279544302$ & $-0.323697314095031872$\\
$C_\eta$ & $2.158752005657785532$ & $0.905729821720199018$\\
$A_\eta$ & $0.164817523964420504$ & $0.217999212918202522$\\
$a_{\eta,1}$ & $-0.307565333551633697$ & $0.230950213208614981$\\
$a_{\eta,2}$ & $0.018726632258564924$ & $-0.004247066120392843$\\
\bottomrule
\end{tabular}
\end{center}

The power series in \eqref{eq:main} is asymptotic in the finite-order
sense. No convergence of its infinite sum is needed or claimed. Numerical
constants in this table are evaluations of the integral and finite
coefficient formulas, rather than fitted parameters.

\section{Exact representations and their signs}

Let $S(k,j)$ denote the Stirling numbers of the second kind, and define
\[
 P_k(t)=\sum_{j=0}^k(-1)^{k-j}S(k,j)t^j,\qquad R_k(t)=\frac{P_k(t)}{t^k}.
\]
Their exponential generating function is
\begin{equation}\label{eq:Pegf}
 \sum_{k\ge0}P_k(t)\frac{z^k}{k!}=\exp(t(1-e^{-z})).
\end{equation}
A proper coloring partitions each independent vertex part into nonempty
monochromatic blocks, and different parts must use different colors.
Thus, writing $q^{\underline j}$ for a falling factorial,
\[
 \chi_{K_\lambda}(q)=\sum_{j_1,\ldots,j_r}
 q^{\underline{j_1+\cdots+j_r}}\prod_iS(\lambda_i,j_i).
\]
Stanley's identity $\AO(G)=(-1)^{|V(G)|}\chi_G(-1)$ \cite{stanley},
followed by the Gamma integral for $(j_1+\cdots+j_r)!$, proves
\begin{equation}\label{eq:exact}
 \AO(K_\lambda)=\int_0^\infty e^{-t}\prod_iP_{\lambda_i}(t)\dd t.
\end{equation}
This is also Sun's exact multipartite integral. Its integrand can have
either sign. Infinite products formed from it are interpreted
coefficientwise, not as globally convergent analytic products.

For a finite part cutoff $K$, put
\[
 F_{\eta,K}(q,t)=\prod_{k=1}^K(1-\eta q^kR_k(t))^{-\eta}.
\]
If $B_{\eta,K}$ denotes the sum restricted to largest part at most $K$,
then, for $T\sim\operatorname{Gamma}(n+1,1)$,
\begin{equation}\label{eq:expectation}
 \frac{B_{\eta,K}(n)}{n!}=\E_T\coeff qn F_{\eta,K}(q,T).
\end{equation}
Coefficient extraction comes first. Every selected profile has total size
$n$, so the removed powers of $t$ contribute exactly $t^n$, producing
the Gamma density in \eqref{eq:expectation}. This finite identity makes
no positivity assumption about the polynomials on the whole half-line.

\section{A uniform quadratic majorant}\label{sec:majorant}

\begin{lemma}\label{lem:roots}
The roots of $P_k$ are distinct and nonnegative, including zero. If
$p_r(k)$ is their $r$th power sum, then it is a polynomial in $k$ of
degree at most $r+1$, determined by
\begin{equation}\label{eq:powerrec}
 p_r(k+1)-p_r(k)=-r\coeff ur
 \log\left(1-ku-\sum_{j\ge1}p_j(k)u^{j+1}\right),\quad p_r(0)=0.
\end{equation}
In particular, the largest root is $O(k^{5/4})$.
\end{lemma}
\begin{proof}
The Stirling recurrence gives
$P_{k+1}(t)=t(P_k(t)-P'_k(t))$, with $P_1(t)=t$.
At consecutive roots of $P_k$, the values of $P_k-P'_k$ alternate in
sign. There is one root in each intervening interval and one to the
right of the largest root; all are positive. Multiplication by $t$
completes the induction.

Writing $\mathcal R_k(u)=\prod_i(1-r_i u)$, we have
$\mathcal R_{k+1}=(1-ku)\mathcal R_k+u^2\mathcal R'_k$.
Taking logarithms yields \eqref{eq:powerrec}. Only $p_j$ with $j<r$
is needed to compute the right side. Its degree in $k$ is at most $r$,
so finite summation proves the degree bound. Explicitly,
\begin{align}
 p_1(k)&=\frac{k(k-1)}2,\qquad
 p_2(k)=\frac{k(k-1)(4k-5)}6,\\
 p_3(k)&=\frac{k(k-1)(9k^2-25k+18)}8,\\
 p_4(k)&=\frac{k(k-1)(64k^3-291k^2+459k-251)}{30}.
 \label{eq:p4}
\end{align}
All roots are nonnegative, so their maximum is at most $p_4(k)^{1/4}$.
\end{proof}

\begin{theorem}\label{thm:majorant}
Fix $3/4<\beta<4/5$. There is a constant $C_0$ such that, uniformly over
partitions $\lambda$ of $n$ whose largest part is at most $n^\beta$,
\begin{equation}\label{eq:majorant}
 \frac{\AO(K_\lambda)}{n!}\le C_0
      \exp\left(-\frac{Q(\lambda)-n}{2n}\right),\qquad
 Q(\lambda)=\sum_i\lambda_i^2.
\end{equation}
The central integral also has Gaussian suppression in
$z=(t-n)/\sqrt n$, uniformly over these profiles.
\end{theorem}
\begin{proof}
Put
\[
 A=\frac{Q-n}{2},\qquad
 B=\sum_i\frac{\lambda_i(\lambda_i-1)(4\lambda_i-5)}{12}.
\]
Lemma~\ref{lem:roots} and $\beta<4/5$ imply that every relevant root is
$o(n)$. All factors are therefore positive on $[n/2,2n]$. The exact
inequality $\log(1-v)\le-v-v^2/2$ gives
\begin{equation}\label{eq:rootmajorant}
 \prod_iP_{\lambda_i}(t)\le t^n\exp(-A/t-B/t^2).
\end{equation}
Weighted Cauchy--Schwarz, followed by a termwise inequality, gives
\[
 \frac{A^2}{n}\le\sum_i\frac{\lambda_i(\lambda_i-1)^2}{4}\le B.
\]
The difference in the second inequality for a part $k$ is
$k(k-1)(k-2)/12\ge0$.

Let $x=t/n$, $a=A/n^2$ and $g(x)=\log x-x+1$. Relative to
$n^ne^{-n-A/n}$, the integrand is bounded by
\[
 \exp\{n[g(x)+a(1-1/x)-a^2/x^2]\}.
\]
For $x\le1$ the two $a$ terms are nonpositive. For $x\ge1$ their
supremum over $a\ge0$ equals $(x-1)^2/4$. On $[1/2,2]$ there is a
constant $\delta>0$ with
\begin{equation}\label{eq:gaussmajorant}
 g(x)+(x-1)^2/4\le-\delta(x-1)^2.
\end{equation}
Indeed its derivative is $(x-1)(1/2-1/x)$; its value at 2 is
$\log2-3/4<0$, and division by $(x-1)^2$ has limiting value $-1/4$
at 1. Integrating the resulting Gaussian bound and using Stirling's
formula proves \eqref{eq:majorant} for the central integral. The same
calculation gives the pointwise factor $e^{-\delta z^2}$.

It remains to control signs outside that interval. Unsigned first-kind
Stirling numbers dominate the second-kind ones: specifying a cyclic
order on each block of a set partition produces a permutation. If
$M=\max_i\lambda_i$, then, for $t\ge0$,
\[
 |P_k(t)|\le\sum_jS(k,j)t^j
 \le t(t+1)\cdots(t+k-1)\le(t+M)^k.
\]
Hence the absolute tail integral is bounded by
\[
 \int_{t\notin[n/2,2n]}e^{-t}(t+M)^n\dd t
      =n!\exp(-\Omega(n)+O(M)).
\]
Substitution of $t+M$ and elementary Gamma Chernoff bounds justify this
uniformly for $M=o(n)$. Since $A/n\le M/2$, the tail is negligible
relative to $n!e^{-A/n}$. Positivity of the orientation count and the
triangle inequality complete the proof.
\end{proof}

\section{Relative precision cutoff removal}\label{sec:cutoff}

Write $s=\sqrt n$ and choose $K=\lfloor s\log n\rfloor$. For the
intermediate range $K<\max\lambda\le n^\beta$, the majorant reduces the
sum to a positive weighted partition model, with
\[
 w_k=e^{-k^2/(2n)},\qquad q_0=e^{-c/s},\qquad
 \mathcal F_\eta(q_0)=\prod_{k\ge1}(1-\eta w_kq_0^k)^{-\eta}.
\]
For the unrestricted model all geometric ratios are below one. In the
distinct model some ratios exceed one, but each factor remains positive
and defines a Bernoulli distribution. A union bound bounds the weighted
coefficient containing a part greater than $K$ by
\begin{equation}\label{eq:union}
 q_0^{-n}\mathcal F_\eta(q_0)
 \begin{cases}
 \displaystyle\sum_{k>K}\frac{w_kq_0^k}{1-w_kq_0^k},&\eta=+1,\\[1ex]
 \displaystyle\sum_{k>K}w_kq_0^k,&\eta=-1.
 \end{cases}
\end{equation}
The sums are bounded by $\exp(-\Omega(\log^2 n))$ times a polynomial.
The Euler--Maclaurin calculation below, already at leading order, gives
$q_0^{-n}\mathcal F_\eta(q_0)=\exp(Cs+O(\log n))$.
Thus this range is negligible relative to $e^{Cs}$ times any fixed
inverse power of $n$.

For the remaining range let $L=\max\lambda>n^\beta$. Splitting every
other part into singleton parts adds edges, so deletion--contraction
implies
\[
 \AO(K_\lambda)\le\AO(K_{L,1,\ldots,1})=(n-L)!(n-L+1)^L.
\]
The formula follows by ordering the singleton vertices and placing each
of the $L$ labeled remaining vertices into one of the $n-L+1$ gaps.
There is an exact bound valid for every $0\le L\le n$:
\begin{align}
 \frac{(n-L)!(n-L+1)^L}{n!}
 &=\prod_{j=1}^L\left(1+\frac{j-1}{n-L+1}\right)^{-1}\\
 &\le\exp\left(-\frac{L(L-1)}{2n}\right).
 \label{eq:fartail}
\end{align}
Use $\log(1+x)\ge x/(1+x)$ and $n-L+j\le n$.
Since $2\beta-1>1/2$, multiplication by the standard bound
$p(n)\le e^{O(\sqrt n)}$ leaves a negligible tail.

The Gamma variable can also be restricted to
\begin{equation}\label{eq:gammawindow}
 |T-n|\le s\log^2 n.
\end{equation}
For profiles with largest part at most $K$, the pointwise form of
\eqref{eq:gaussmajorant} gives suppression $e^{-\delta z^2}$ on
$[n/2,2n]$, in addition to the weight $e^{-(Q-n)/(2n)}$. Summing these
fixed-total weights gives the same positive weighted-partition
coefficient as above. The complement of \eqref{eq:gammawindow} is
therefore smaller than every algebraic order on the main scale.
Outside $[n/2,2n]$ the absolute signed-tail estimate applies.
This argument uses positivity only where it was proved.

\section{The local product and Euler--Maclaurin expansion}\label{sec:EM}

Set $\eps=1/s$ and $T=s^2+sz$. On the cutoffs already justified,
\begin{equation}\label{eq:rootseries}
 \log R_k(T)=-\sum_{r\ge1}\frac{p_r(k)}{rT^r}.
\end{equation}
This convergent root expansion has a useful finite remainder bound:
after $r=R$, its absolute remainder is at most
\begin{equation}\label{eq:rooterror}
 \frac{p_{R+1}(k)}{(R+1)T^{R+1}(1-r_{\max}/T)}.
\end{equation}
Here $r_{\max}/T=o(1)$ uniformly and
$p_{R+1}(k)=O_R((1+k)^{R+2})$. Consequently, with $x=k\eps$,
\[
 \log R_k(s^2+sz)=-x^2/2+\sum_{j\ge1}\eps^j\ell_j(x,z)
\]
is a controlled arbitrary finite expansion. Each $\ell_j$ is a
polynomial vanishing at $x=0$. The first two are
\begin{align}
 \ell_1&=x/2+zx^2/2-x^3/3,\label{eq:ell1}\\
 \ell_2&=-zx/2-z^2x^2/2+2zx^3/3+3x^2/4-3x^4/8.
 \label{eq:ell2}
\end{align}

Define $f_j$ by formal coefficient extraction in
\begin{equation}\label{eq:fj}
 -\eta\log\left(1-\eta\exp\left[-cx-x^2/2+
                  \sum_{j\ge1}\eps^j\ell_j(x,z)\right]\right)
 =f_0(x;c)+\sum_{j\ge1}\eps^jf_j(x;c,z).
\end{equation}
For $j\ge1$, $f_j$ is smooth at zero, including in the Bose case.
The $m$th derivative of the outer logarithm has a pole of order at most
$m$, whereas a product of $m$ perturbation factors vanishes to at least
order $m$. At infinity all relevant derivatives decay as a Gaussian
times a polynomial. These properties persist for $c$ in a sufficiently
small complex neighborhood of its real saddle, with polynomial bounds
in $z$.

In the Fermi case $f_0$ is smooth and $f_0(0)=\log2$. In the Bose case
subtract $f_{\rm std}(x;c)=-\log(1-e^{-cx})$ and put
$d_0(x;c)=f_0-f_{\rm std}$. Then $d_0$ is smooth with $d_0(0)=0$.
The standard product expansion is
\begin{equation}\label{eq:eta}
 \sum_{k\ge1}f_{\rm std}(k\eps;c)
 =\frac{\pi^2}{6c\eps}+\frac12\log\frac{c\eps}{2\pi}
       -\frac{c\eps}{24}+O_R(\eps^R)
\end{equation}
for every fixed $R$. One proof uses Mellin inversion with kernel
$\Gamma(v)\zeta(v)\zeta(v+1)$. Apart from $v=1,0,-1$, the poles at
negative integers cancel by the trivial zeta zeros; moving the contour
to any fixed negative line proves the stated finite-order version.
Only this version is required here.

Euler--Maclaurin now gives
\begin{equation}\label{eq:productexp}
 \log F_{\eta,K}(e^{-c\eps},T)
 =\frac{J(c)}\eps+r_\eta\log\eps+d_\eta(c)
          +\sum_{j\ge0}\eps^jH_j(c,z),
\end{equation}
where
\begin{align}
 r_+&=\tfrac12,&d_+(c)&=\tfrac12\log(c/(2\pi)),\\
 r_-&=0,&d_-(c)&=-\tfrac12\log2,\\
 H_0(c,z)&=\frac{M_1(c)}2+\frac{zM_2(c)}2-\frac{M_3(c)}3.
 \label{eq:H0}
\end{align}
For $j\ge1$ the explicit coefficient formula is
\begin{equation}\label{eq:Hj}
 H_j=\int_0^\infty f_{j+1}\dd x-\frac{f_j(0)}2
 -\sum_{m=1}^{\lfloor j/2\rfloor}\frac{\mathsf B_{2m}}{(2m)!}
        f_{j-2m+1}^{(2m-1)}(0)+E_j(c),
\end{equation}
where $\mathsf B_j$ are Bernoulli numbers, and superscripts on $f_j$
in \eqref{eq:Hj} mean $x$ derivatives.
For the Fermi case,
\[
 E_{2m}=0,\qquad E_j=-\frac{\mathsf B_{j+1}}{(j+1)!}f_0^{(j)}(0)
 \quad(j\text{ odd}).
\]
For the Bose case,
\[
 E_{2m}=0,\qquad E_1=\frac{1/c-c}{24},\qquad
 E_j=-\frac{\mathsf B_{j+1}}{(j+1)!}d_0^{(j)}(0)
 \quad(j\ge3\text{ odd}).
\]
Thus only finitely many polynomials, integrals and endpoint derivatives
are needed for any prescribed order. The formal coefficient functions
can be extended to $k>K$ because their Gaussian tails are smaller than
every algebraic order; this does not extend the original signed infinite
product as an analytic object.

\section{Fourier localization and the prefactors}\label{sec:fourier}

For a central value of $T$, every $P_k(T)$ in the cutoff product is
positive. On $q=q_0e^{\ii\theta}$ the product ratio is a characteristic
function of independent geometric or Bernoulli variables. Choose a fixed
interval $as\le k\le bs$ with $0<a<b$. Its parameters stay uniformly
away from degeneracy, so for some positive constants $d,d'$,
\begin{align}
 \frac{|F(q_0e^{\ii\theta},T)|}{F(q_0,T)}
 &\le\exp\left(-d\sum_{as\le k\le bs}(1-\cos(k\theta))\right)\\
 &\le\exp(-d'\min(s^3\theta^2,s)),\qquad |\theta|\le\pi.
 \label{eq:arcs}
\end{align}
For $|\theta|\le c_0/s$, use the quadratic bound on $1-\cos(k\theta)$.
For $c_0/s\le|\theta|\le C_0/s$, use compactness and the nonzero
limiting integral of $1-\cos(xv)$. For $|\theta|\ge C_0/s$, the
geometric-sum estimate for $\sum e^{\ii k\theta}$ gives a positive
fraction of the interval length when $C_0$ is large enough. This covers
the full circle, including possible rational angles.

Only $|\theta|\le s^{-3/2}\log^2 n$ contributes at algebraic order.
Write $\theta=us^{-3/2}$. In \eqref{eq:productexp} this replaces $c$ by
$c'=c-\ii u\sqrt\eps$. The linear phase from $q^{-n}$ cancels because
$J'(c)=-M_1(c)=-1$, and the quadratic phase is $-Vu^2/2$.

The density of $z=(T-n)/s$ has the finite expansion
\begin{align}\label{eq:gammadensity}
 \frac{e^{-z^2/2}}{\sqrt{2\pi}}
 \exp\bigg(&\sum_{j\ge1}\frac{(-1)^{j+1}}{j+2}\eps^jz^{j+2}
 \\
 &-\sum_{m\ge1}\frac{\mathsf B_{2m}}{2m(2m-1)}\eps^{4m-2}\bigg).
 \nonumber
\end{align}
It is obtained directly from the Gamma density and Stirling's formula;
in particular its first correction is $\eps z^3/3$, with no linear
$z$ term when centering at $n$.

At the saddle, the constant product amplitude is
\[
 e^{H_0}=\exp(1/2-M_3/3)\exp(zM_2/2).
\]
Integration against the normal density contributes
$\exp(1/2-M_3/3+M_2^2/8)=E_0$.
The Fourier integral contributes $(2\pi s^3V)^{-1/2}$. The
Euler--Maclaurin boundary contributes $\sqrt{c/(2\pi s)}$ for Bose
and $1/\sqrt2$ for Fermi. Their product is exactly \eqref{eq:A},
with the powers $n^{-1}$ and $n^{-3/4}$. Omitting the Gamma tilt or
the cubic Stirling-transform contribution would change the amplitudes.

\section{A finite coefficient algorithm at every order}\label{sec:generator}

Let $U$ and $Z$ be independent Gaussian variables with
\[
 U\sim N(0,V^{-1}),\qquad Z\sim N(\alpha,1),\qquad\alpha=M_2/2.
\]
Use $\sqrt\eps$ as a formal variable and set $\Delta=-\ii U\sqrt\eps$.
Define the formal exponent
\begin{align}
 \mathcal Q(\eps,U,Z)={}&
 \sum_{r\ge3}\frac{J^{(r)}(c)}{r!}(-\ii U)^r\eps^{r/2-1}
 +d(c+\Delta)-d(c)\nonumber\\
 &+H_0(c+\Delta,Z)-H_0(c,Z)
 +\sum_{j\ge1}\eps^jH_j(c+\Delta,Z)\nonumber\\
 &+\sum_{j\ge1}\frac{(-1)^{j+1}}{j+2}\eps^jZ^{j+2}
 -\sum_{m\ge1}\frac{\mathsf B_{2m}}{2m(2m-1)}\eps^{4m-2}.
 \label{eq:masterQ}
\end{align}
The complete coefficient prescription is
\begin{equation}\label{eq:master}
 \boxed{\displaystyle a_j=\coeff\eps j\E\exp(\mathcal Q(\eps,U,Z)).}
\end{equation}
For fixed $j$, this requires only finite truncations of
\eqref{eq:powerrec}, \eqref{eq:fj}, \eqref{eq:Hj} and
\eqref{eq:masterQ}, followed by Gaussian moments. Odd half-integer
powers integrate to zero because their integrands are odd in $U$.
This proves that the expansion scale is $n^{-1/2}$ with no extra
logarithmic factors.

\subsection{Uniform remainder control}

The cutoffs in largest part and Gamma variable have already been removed
at relative precision in Section~\ref{sec:cutoff}. Bound \eqref{eq:arcs}
is uniform in the remaining Gamma window. In the local product, the
positive-root remainder \eqref{eq:rooterror} allows arbitrarily many
terms to be retained. After removing the Bose logarithmic endpoint,
Euler--Maclaurin remainders are controlled by integrals of derivatives;
these are finite and at most polynomial in $|z|$, uniformly in a fixed
complex neighborhood of $c$.

A further detail is important. The fixed-total profile majorant by
itself does not bound the integrand $F(q_0,T)$, which includes all total
sizes. Its joint Gaussian envelope follows instead from the local
product expansion. Uniformly for $|z|\le\log^2 n$,
\begin{equation}\label{eq:envelopeproduct}
 \log F(q_0,T)=sJ+r_\eta\log\eps+d_\eta+H_0(c,z)
                 +O(\eps(1+z^2)).
\end{equation}
For completeness, \eqref{eq:rootseries} gives on the cutoff range
\[
 \log R_k+x^2/2=\eps\ell_1+O(\eps^2(1+z^2)P(x)),\qquad P(0)=0,
\]
for a fixed polynomial $P$. The endpoint factor follows from the
vanishing of the root power-sum polynomials at $k=0$. The outer
logarithm's linear derivative times $P$ and quadratic derivative times
$\ell_1^2$ are integrable even in the Bose case. Summation over the
mesh multiplies this bound by $s$; Gaussian decay in $x$ controls all
polynomial factors. Endpoint corrections have the same bound, proving
\eqref{eq:envelopeproduct}.

Because $H_0$ is affine in $z$, multiplication by the exact Gamma
density gives, after extracting common scale factors, an upper bound
\[
 C_1\exp\{-z^2/2+\alpha z+O(\eps(1+|z|^3))\}
          \le C_2e^{-z^2/4}.
\]
The error is absorbed since $\eps\log^2 n\to0$. Together with
\eqref{eq:arcs}, this supplies a joint Gaussian envelope in $u,z$.
Finite Taylor remainders can therefore be bounded by a fixed polynomial
in $|u|+|z|$ times the target power of $\eps$, multiplied by
$e^{-d_2u^2-d_3z^2}$ for positive constants $d_2,d_3$.
The resulting integrals are bounded independently of $n$.
Expanding a few extra orders if necessary and integrating the next odd
half-power, which vanishes exactly, gives an $O_L(\eps^{L+1})$
remainder. No logarithmic cutoff losses remain. This justifies
\eqref{eq:master} at every fixed order and completes the proof of
Theorem~\ref{thm:main}.

\section{First corrections and executable calculations}

Define the cumulant integrals
\[
 K_j=\int_0^\infty x^j\rho(1+\eta\rho)\dd x.
\]
Then the first product correction is the quadratic polynomial
\begin{align}
 H_1(z)={}&e_\eta+\frac{3M_2}{4}-\frac{3M_4}{8}
            +\frac{K_2}{8}-\frac{K_4}{6}+\frac{K_6}{18}\nonumber\\
 &+z\left(-\frac{M_1}{2}+\frac{2M_3}{3}
                         +\frac{K_3}{4}-\frac{K_5}{6}\right)
 +z^2\left(-\frac{M_2}{2}+\frac{K_4}{8}\right),\label{eq:H1}\\
 e_+={}&-\frac c{24}-\frac5{24c},\qquad e_- =\frac c{24}.
 \nonumber
\end{align}
Let $A_0(c,z)=e^{d(c)+H_0(c,z)}$, with $c$ left variable while
differentiating. The Fourier correction is
\begin{equation}\label{eq:D}
 D(z)=-\frac{(A_0)_{cc}}{2VA_0}
 +\frac{J'''(A_0)_c}{2V^2A_0}
 +\frac{J''''}{8V^2}-\frac{5(J''')^2}{24V^3}.
\end{equation}
Consequently
\begin{equation}\label{eq:a1}
 a_1=\E_{Z\sim N(\alpha,1)}\left[H_1(Z)+Z^3/3+D(Z)\right].
\end{equation}
This provides an independent short route to the first coefficient.
The script \texttt{first\_correction.py} evaluates it by
one-dimensional integrals and exact Gaussian moments.

The separate script \texttt{generic\_second\_correction.py} implements
\eqref{eq:master} through $a_2$ without using the hand-expanded
\eqref{eq:H1} or \eqref{eq:D}. It uses the displayed root polynomials
to construct the $\ell_j$, cumulants, and all needed $c$ derivatives. It
reproduces both $a_1$ values and gives the $a_2$ values in the table.
Higher orders use the same finite algorithm; the supplied executable
is tested through second order rather than advertised as a fully
optimized arbitrary-order software package.

\section{Controlled smooth inverses}\label{sec:inverse}

The sequence itself has only integer arguments. We first specify the
sense of a smooth inverse. Let $G(x)$ be a smooth realization of its
full logarithmic asymptotic expansion, with termwise differentiable
asymptotics. Such a realization is obtained by multiplying successive
terms by fixed smooth cutoffs at sufficiently rapidly increasing
radii, chosen inductively to control finitely many additional
derivatives at each step.

If exact interpolation is desired, choose a smooth bump $\phi$ supported
in $(-1/3,1/3)$ with $\phi(0)=1$ and define
\begin{equation}\label{eq:interpolant}
 \log F(x)=G(x)+\sum_{\substack{m\text{ integer}\\m\text{ sufficiently large}}}
       (\log B(m)-G(m))\phi(x-m).
\end{equation}
The supports are disjoint, and the correction coefficients and all
their bump derivatives are smaller than every inverse power by
Theorem~\ref{thm:main}. Thus $F(m)=B(m)$, $F$ is positive, and
$(\log F)'(x)=\log x+O(x^{-1/2})$. It is eventually strictly increasing.
All such realizations have the same inverse expansion to every
displayed order. No canonical analytic continuation is asserted.

Let
\[
 \log\left(1+\sum_{j\ge1}a_jx^{-j/2}\right)
       =\sum_{j\ge1}b_jx^{-j/2}.
\]
Stirling's formula gives
\begin{equation}\label{eq:logmodel}
 \log F(x)=x(\log x-1)+C\sqrt x+\alpha_0\log x+d_0
                         +\sum_{j\ge1}\beta_jx^{-j/2},
\end{equation}
where $\alpha_0=1/2-p$, $d_0=\log(A\sqrt{2\pi})$,
$\beta_1=a_1$, and $\beta_2=b_2+1/12$.
For $Y\to\infty$, put
\[
 L=\log Y,\qquad N=\frac{L}{W(L/e)},\qquad\ell=\log N,
\]
with the positive real branch of Lambert $W$.

\begin{theorem}\label{thm:inverse}
The smooth inverse admits a finite-order expansion of the form
\[
 F^{-1}(Y)=N+\sum_{j\ge-1}N^{-j/2}r_j(1/\ell),
\]
where each $r_j$ is a polynomial in $1/\ell$, computed recursively
from \eqref{eq:logmodel}. In particular,
\begin{align}\label{eq:inverse}
 F^{-1}(Y)={}&N-\frac{C\sqrt N}{\ell}-\alpha_0-\frac{d_0}{\ell}
 +\frac{C^2}{2\ell^2}-\frac{C^2}{2\ell^3}
 +O\left(\frac1{\sqrt N\ell}\right).
\end{align}
The constant shifts are $+1/2$ for unrestricted partitions and
$+1/4$ for distinct parts.
\end{theorem}
\begin{proof}
Substitute $x=N+\Delta$ in \eqref{eq:logmodel} and use
$N(\log N-1)=L$. At each new power of $N^{-1/2}$, the unknown
coefficient occurs multiplied by $\ell$, so it is determined uniquely
by division by $\ell$. Taylor expansion first cancels the
$\sqrt N$ term, then the constant term, yielding \eqref{eq:inverse}.
All remaining operations preserve polynomials in $1/\ell$.

For quantitative control, substitute a finite inverse truncation $X$
into the logarithmic expansion and compute its residual to one more
order. If $\log F(X)-L=O(N^{-r/2})$, the mean-value theorem and
$(\log F)'\sim\log N$ give inverse error
$O(N^{-r/2}/\log N)$. This proves each prescribed finite order;
it does not invert only the leading equivalent.
\end{proof}

A useful next term can be obtained without a large expanded polynomial.
Define
\begin{gather}
 a=-C/\ell,\qquad
 b=-\alpha_0-d_0/\ell+C^2/(2\ell^2)-C^2/(2\ell^3),\nonumber\\
 d=-\frac{(a+C/2)b-a^3/6-Ca^2/8+\alpha_0a+\beta_1}{\ell}.
 \label{eq:inversenext}
\end{gather}
Then $N+\sqrt N a+b+d/\sqrt N$ has error $O(1/(N\ell))$.

For the integer threshold $\tau(Y)=\min\{n:B(n)\ge Y\}$,
eventual monotonicity follows from the main expansion. With an exact
increasing interpolant as in \eqref{eq:interpolant},
$\tau(Y)=\lceil F^{-1}(Y)\rceil$ for sufficiently large $Y$.
A finite approximation $X(Y)$ with a valid error bound $e(Y)$ gives
only the rigorous envelope
\begin{equation}\label{eq:threshold}
 \lceil X(Y)-e(Y)\rceil\le\tau(Y)\le\lceil X(Y)+e(Y)\rceil.
\end{equation}
Exact rounding follows when the endpoints coincide, not unconditionally
from a finite asymptotic formula.

\section{Exact checks and reproducibility}

The program \texttt{exact\_check.py} computes both sequences by positive
Touchard polynomial arithmetic followed by one exact alternating
factorial evaluation. Let $Q_k(t)=\sum_jS(k,j)t^j$. Compute
\[
 H_{+,n}(t)=\coeff qn\prod_{k=1}^n(1-q^kQ_k(t))^{-1},\qquad
 H_{-,n}(t)=\coeff qn\prod_{k=1}^n(1+q^kQ_k(t)).
\]
If $h_{n,j}$ are their coefficients, then
$B_\eta(n)=\sum_j(-1)^{n-j}h_{n,j}j!$.
All intermediate $h_{n,j}$ are nonnegative. For a fixed profile the
sum of its coefficients is at most the Bell number $\operatorname{Bell}(n)$:
its blockwise set partitions inject into set partitions of a fixed
$n$-set. Hence
\[
 h_{n,j}\le p(n)\operatorname{Bell}(n)\le2^n n^n.
\]
The binary base in the Kronecker encoding is chosen above this bound,
so no polynomial digit can carry. All sequence values are exact Python
integers; cancellation occurs only in the final exact evaluation.

Separate runs at maximum sizes 150 and 300 agree on every overlapping
value, and all 21 initial OEIS terms used in the check agree. Let
\[
 R_\eta(n)=\frac{B_\eta(n)}{A_\eta n!n^{-p_\eta}e^{C_\eta\sqrt n}},\qquad
 T_\eta(n)=n^{3/2}\left(R_\eta(n)-1-\frac{a_{\eta,1}}{\sqrt n}
                                      -\frac{a_{\eta,2}}n\right).
\]
Representative exact-data residuals are
\begin{center}
\begin{tabular}{@{}rrrrr@{}}
\toprule
$n$ & $R_+(n)$ & $T_+(n)$ & $R_-(n)$ & $T_-(n)$\\
\midrule
100 & 0.9694277210 & -0.00301198 & 1.0230946532 & 0.04210253\\
150 & 0.9750093079 & -0.00538650 & 1.0188567728 & 0.05158749\\
200 & 0.9783435022 & -0.00559355 & 1.0163242684 & 0.04202347\\
250 & 0.9806213569 & -0.00557333 & 1.0146001519 & 0.04176652\\
300 & 0.9823040643 & -0.00553406 & 1.0133277008 & 0.04126194\\
\bottomrule
\end{tabular}
\end{center}
The slower and mildly oscillatory finite-size behavior in the distinct
case does not enter the proof. These checks support, but do not replace,
the analytic argument.

The verification script also tests the inverse at $Y=B_\eta(300)$.
Formula \eqref{eq:inverse} differs from 300 by approximately
$0.00350453$ and $0.00380904$, respectively. Including
\eqref{eq:inversenext} changes these to approximately
$-0.0000155943$ and $0.0000217153$.
These are numerical residual checks, not certified threshold bounds.

The reproducibility package contains the TeX source, PDF, proof source,
exact arrays, numerical outputs, verification scripts and review record.
Its \texttt{README.md} specifies the short replay and full exact replay.
The second-order executable uses Python with \texttt{sympy} and
\texttt{mpmath}; no proprietary algebra system is required.

\section{A fixed positive chromatic parameter}

The proof also treats the natural chromatic specialization
$H_v(G)=(-1)^{|V(G)|}\chi_G(-v)$ for fixed $v>0$.
Let $B_{\eta,v}(n)$ sum $H_v(K_\lambda)$ over the same two partition
families. This section asserts uniformity only when $v$ ranges over a
fixed compact subset of $(0,\infty)$.

\begin{corollary}\label{cor:chromatic}
The constants $A_\eta,C_\eta,p_\eta$ in Theorem~\ref{thm:main} are
unchanged, and
\begin{equation}\label{eq:chromatic}
 B_{\eta,v}(n)=\frac{\Gamma(n+v)}{\Gamma(v)}
 A_\eta n^{-p_\eta}e^{C_\eta\sqrt n}
 \left(\sum_{j=0}^L a_{\eta,j}(v)n^{-j/2}
                   +O_L(n^{-(L+1)/2})\right).
\end{equation}
The coefficients have the same finite construction, and
\[
 a_{\eta,1}(v)=a_{\eta,1}(1)+(v-1)M_2/2.
\]
\end{corollary}
\begin{proof}
Replace \eqref{eq:exact} by the exact identity
\[
 H_v(K_\lambda)=\frac1{\Gamma(v)}\int_0^\infty
 e^{-t}t^{v-1}\prod_iP_{\lambda_i}(t)\dd t.
\]
The normalized expectation uses $T\sim\operatorname{Gamma}(n+v,1)$.
In the formal exponent \eqref{eq:masterQ}, replace the Gamma part by
\begin{align*}
 &\sum_{j\ge1}\frac{(-1)^{j+1}}{j+2}\eps^jZ^{j+2}
 +(v-1)\log(1+\eps Z)\\
 &\hspace{1cm}
 -\sum_{r\ge1}\frac{(-1)^{r+1}\mathsf B_{r+1}(v)}{r(r+1)}\eps^{2r},
\end{align*}
where $\mathsf B_r(v)$ is a Bernoulli polynomial. This follows from
generalized Stirling expansion of $\Gamma(n+v)$. The only first-order
change is $(v-1)Z$, with tilted mean $(v-1)M_2/2$.

On the central interval $t/n\in[1/2,2]$, the extra factor $t^{v-1}$
is uniformly comparable to $n^{v-1}$, and
$\Gamma(n+v)\asymp n!n^{v-1}$. All majorants remain uniform. For the
lower signed tail, including $0<v<1$, do not bound $t^{v-1}$ pointwise
near zero. Instead $e^{-t}(t+M)^n$ is increasing on $[0,n/2]$ for large
$n$, and its endpoint bound is multiplied by
$\int_0^{n/2}t^{v-1}\dd t=(n/2)^v/v$. This gives the same exponential
suppression uniformly for $v$ bounded away from zero. The upper tail
is unchanged by fixed powers.

Finally $H_v$ is positive: chromatic coefficients have alternating
nonnegative signs. Deletion--contraction gives
$H_v(G+e)=H_v(G)+H_v(G/e)$, hence edge monotonicity. With $m=n-L$,
\[
 H_v(K_{L,1,\ldots,1})=\frac{\Gamma(m+v)}{\Gamma(v)}(m+v)^L.
\]
Dividing by $\Gamma(n+v)/\Gamma(v)$ gives, by the same logarithmic
inequality as \eqref{eq:fartail}, the bound
$\exp(-L(L-1)/(2(n+v-1)))$. The far-part tail and every local estimate
therefore transfer to compact positive $v$ sets.
\end{proof}

For the inverse, normalize \eqref{eq:chromatic} relative to
$n!n^{v-1}$. In \eqref{eq:logmodel} this changes the constants to
\[
 \alpha_0=v-\tfrac12-p_\eta,\qquad
 d_0=\log\left(\frac{A_\eta\sqrt{2\pi}}{\Gamma(v)}\right).
\]
The same Lambert-$W$ core and recursion apply. The Gamma-ratio series
affects $\beta_2$ and higher terms, but not
$\beta_1=a_{\eta,1}(v)$.

\section{Further questions}

The complete algebraic expansion leaves several substantial questions.
First, one can study the late growth of its coefficients and the
exponentially small sectors arising from noncentral coefficient
contours; the present theorem does not identify these sectors or a
summation prescription. Second, the dominant weighted partition model
suggests limit laws for the largest part and for profile statistics
under orientation weighting. Turning this into uniform distributional
results requires additional marked-parameter estimates. Third, a
chromatic parameter growing with $n$ could change the Gamma tilt and
the saddle regime; compact-parameter arguments do not cover that
transition. Finally, explicit effective constants in the full
remainder would turn the inverse threshold envelope into a practical
certified numerical algorithm.

\begin{thebibliography}{9}
\bibitem{oeisb} The OEIS Foundation,
\href{https://oeis.org/A372395}{A372395}, total acyclic orientations over
complete multipartite graphs indexed by integer partitions.
Accessed 1 October 2026.
\bibitem{oeisf} The OEIS Foundation,
\href{https://oeis.org/A370613}{A370613}, the corresponding distinct-part
partition sum. Accessed 1 October 2026.
\bibitem{stanley} R.~P.~Stanley, Acyclic orientations of graphs,
\emph{Discrete Mathematics} \textbf{5} (1973), 171--178.
\href{https://doi.org/10.1016/0012-365X(73)90108-8}{doi:10.1016/0012-365X(73)90108-8}.
\bibitem{sun} Z.~Sun, Saddle-Point Asymptotics for Chromatic and Tutte
Polynomial Evaluations of Complete Multipartite Graphs,
\href{https://arxiv.org/abs/2605.04006v1}{arXiv:2605.04006v1} (2026),
especially Sections 2, 7--9 and 11.
\end{thebibliography}
\end{document}
